\chapter{Teoría de grafos y recocido simulado}
\section{Fundamentos y definiciones formales}
Un grafo simple no dirigido es un par $G=(V,E)$, donde $V$ es un conjunto finito de vértices y $E$ es un conjunto de pares no ordenados de vértices distintos. No hay lazos ni aristas múltiples. Se trabajará con $n=|V|\geq3$ y $m=|E|$. Una arista representa una conexión, no necesariamente una distancia física. Un recorrido es una secuencia $(v_0,\ldots,v_k)$ tal que $\{v_{i-1},v_i\}\in E$ para $1\leq i\leq k$; su longitud es $k$, el número de aristas recorridas. Un camino simple es un recorrido cuyos vértices son distintos. Un ciclo simple es una secuencia $(v_0,\ldots,v_{k-1},v_0)$, con $k\geq3$, cuyos primeros $k$ vértices son distintos y todas las conexiones existen. Se explicita ``simple'' porque las convenciones terminológicas varían entre textos.\cite{sa1,sa2}

Un camino hamiltoniano contiene todos los vértices exactamente una vez y tiene $n-1$ aristas. Un ciclo hamiltoniano contiene todos los vértices y tiene longitud exactamente $n$. Formalmente, para una permutación $\pi=(v_1,\ldots,v_n)$ de $V$, se requiere
\[
\{v_i,v_{i+1}\}\in E\quad(1\leq i<n),\qquad \{v_n,v_1\}\in E.
\]
El primer vértice puede escribirse nuevamente al final para indicar el cierre; no es una segunda visita intermedia. No es necesario recorrer todas las aristas del grafo. El TSP añade costos y busca minimizar el costo del recorrido; aquí se busca existencia, por lo que todos los ciclos válidos tienen la misma longitud combinatoria $n$.\cite{sa2,sa3}

\section{Verificación de una ruta}
Para comprobar una candidata se verifica que contiene exactamente $n$ identificadores válidos, que no se repite ninguno, que aparecen todos los vértices y que existen las $n$ conexiones, incluida la última hacia el origen. Con consultas de adyacencia de costo constante, esta verificación cuesta $O(n)$ y requiere $O(n)$ espacio auxiliar si se usa una marca por vértice. La complejidad de encontrar una solución no debe confundirse con la de comprobarla.

Considérese $V=\{A,B,C,D,E\}$ y $E=\{AB,BC,CD,DE,EA\}$. La ruta $A,B,C,D,E,A$ es un ciclo de cinco aristas. La permutación $A,B,D,C,E$ contiene todos los vértices, pero falla en $BD$ y $CE$. Si se elimina $EA$ del grafo, $A,B,C,D,E$ sigue siendo un camino hamiltoniano, pero deja de poder cerrarse. Estos ejemplos proceden de la aportación de grafos y SA; son didácticos.

\section{Fundamentos del recocido simulado}
SA se inspira en el enfriamiento de materiales. Mantiene un estado actual, propone un vecino y compara sus costos. La temperatura $T>0$ controla la aceptación de empeoramientos; no representa grados físicos en el programa. Para minimizar $b$, se define $\Delta=b(\pi')-b(\pi)$ y se acepta la candidata con probabilidad
\[
P(\text{aceptar})=\begin{cases}1,&\Delta\leq0,\\ \exp(-\Delta/T),&\Delta>0.\end{cases}
\]
La comparación se realiza con el estado actual y se conserva por separado el mejor observado. Aceptar empeoramientos puede permitir escapar de mínimos locales. Al bajar $T$, un mismo empeoramiento se acepta con menor probabilidad; las mejoras siguen aceptándose.\cite{sa4,sa8}

Hajek estudia condiciones de convergencia con calendarios de enfriamiento específicos. Esas garantías asintóticas no implican que una ejecución finita con enfriamiento geométrico encuentre siempre un ciclo.\cite{sa5}

\section{Representación, costo y vecindario propuestos}
Se representa cada estado como una permutación de $V$, con el regreso implícito. Fijar un vértice al comienzo elimina duplicados por rotación sin perder ciclos. Se propone minimizar
\[
b(\pi)=\sum_{i=1}^{n}\mathbf{1}[\{v_i,v_{i+1}\}\notin E],\qquad v_{n+1}=v_1.
\]
La función cuenta conexiones ausentes: no añade aristas al grafo. Sobre permutaciones, $b=0$ equivale a un ciclo hamiltoniano. Un valor positivo mide incumplimiento y no certifica una solución. No conviene usar la longitud como segundo objetivo: una permutación circular ya tiene exactamente $n$ conexiones candidatas.

Los movimientos propuestos son intercambiar dos posiciones o invertir un segmento, sin mover el inicio fijo. Ambos conservan los vértices, pero pueden romper conexiones. En grafos dirigidos, invertir un segmento también cambia el sentido de sus aristas internas; esta adaptación se limita a grafos no dirigidos. Restringir la búsqueda a ciclos válidos exigiría disponer de uno inicialmente, lo cual ya resolvería el problema de existencia.

\section{Ejemplo didáctico y parámetros}
En el grafo de cinco vértices anterior se parte de $A,B,D,C,E$, con $b=2$. Se elige $T_0=2$, un factor de enfriamiento $\alpha=0.8$ y un intento por temperatura, únicamente para ilustrar el método.
\begin{enumerate}
\item Intercambiar $B$ y $D$ produce $A,D,B,C,E$, con $b=3$: faltan $AD$, $DB$ y $CE$. Como $\Delta=1$, la aceptación tiene probabilidad $e^{-1/2}\approx0.607$. Un valor ilustrativo $u=0.40$ produce aceptación. El mejor costo guardado sigue siendo 2 y la temperatura baja a 1.6.
\item Invertir $D,B,C$ produce $A,C,B,D,E$, con $b=2$. Faltan $AC$ y $BD$; se acepta la mejora y la temperatura baja a 1.28.
\item Intercambiar $C$ y $B$ produce $A,B,C,D,E$, con $b=0$. Las aristas $AB,BC,CD,DE,EA$ existen, por lo que se verifica el ciclo y termina la búsqueda.
\end{enumerate}
Los movimientos y el azar se escogieron a mano; no constituyen una medición ni indican que SA deba resolver FHCP en tres pasos.

\section{Pseudocódigo de la adaptación}
\begin{enumerate}
\item Recibir $G$, temperaturas $T_0>T_{\min}>0$, $0<\alpha<1$, intentos por nivel $L\geq1$ y presupuesto total $B\geq1$.
\item Crear una permutación inicial, evaluar su costo y conservarla como actual y mejor. Verificarla y terminar si ya es un ciclo. Contar también esta evaluación inicial.
\item Mientras queden recursos y $T>T_{\min}$, realizar hasta $L$ propuestas: elegir dos posiciones distintas y aplicar intercambio o inversión con probabilidad $1/2$.
\item Antes de evaluar cada propuesta, comprobar el presupuesto. Evaluarla y aumentar el contador, tanto si se acepta como si se rechaza.
\item Aceptar si $\Delta\leq0$; en otro caso, generar $u$ uniforme en $[0,1)$ y aceptar si $u<e^{-\Delta/T}$.
\item Actualizar el mejor estado y comprobar de forma independiente cualquier candidata con costo cero. Si pasa la verificación, devolverla.
\item Tras cada nivel, actualizar $T\leftarrow\alpha T$. Si no se ha detenido, repetir.
\item Al terminar sin ciclo, devolver la mejor candidata y el estado ``no encontrado dentro del presupuesto''.
\end{enumerate}
El esquema adapta las operaciones descritas en la literatura de SA.\cite{sa4,sa7} Temperatura, calendario, vecindario y escala del costo deben ajustarse conjuntamente. Un enfriamiento rápido reduce exploración; uno lento exige más recursos. Las mesetas y los mínimos locales siguen siendo posibles. Un fallo no demuestra que el grafo sea no hamiltoniano.
